\vspace{2.5pt plus 3pt minus 0.5pt}\noindent\hspace*{1.2em}\kbvnum{೨.}\hspace{0.55em}ತಿನ್ನುತ್ತಿದೆ ೪. ಕರಿಯುತ್ತಿದ್ದಾಳೆ\par

\vspace{2.5pt plus 3pt minus 0.5pt}\noindent\hspace*{1.2em}\kbvnum{೩.}\hspace{0.55em}ಮಾಡುತ್ತಿದ್ದಾನೆ ೫. ಆಡುತ್ತಿದ್ದಾರೆ.\par

\kbvchapter{ಪಾಠ ೫}{ಕ್ರಿಯಾಪದ}

\indent ಕ್ರಿಯೆಯ ಅರ್ಥವನ್ನು ಕೊಡುವ ಪದವೇ ಕ್ರಿಯಾಪದ; ಇದು ೪ ನಿರ್ವಹಿಸುವ ಕೆಲಸವನ್ನು ತಿಳಿಸುತ್ತದೆ.\par

\begin{kbvderiv}
\vspace{4pt plus 3pt minus 1pt}
\noindent{\renewcommand{\arraystretch}{1.24}\setlength{\tabcolsep}{5pt}
\begin{tabularx}{\textwidth}{X X X}
ಕುದುರೆಯು ಓಡುತ್ತಿದೆ & ಪದ್ಮಿನಿಯು ಹಾಡುತ್ತಾಳೆ & ನಾಯಿ \\
ರೊಟ್ಟಿಯನ್ನು ತಿನ್ನುತ್ತದೆ & ಮಂಜುನಾಥನು ಕಾಲ್ಚೆಂಡಿನ ಆಟ &  \\
\end{tabularx}}
\vspace{4pt plus 3pt minus 1pt}
\end{kbvderiv}

\indent ಆಡುತ್ತಾನೆ--ಈ ವಾಕ್ಯಗಳಲ್ಲಿರುವ ಕ್ರಿಯಾಪದಗಳೆಂದರೆ ಓಡು, ಹಾಡುತ್ತಾಳೆ, ತಿನ್ನುತ್ತದೆ, ಆಡುತ್ತಾನೆ. ಈ ಕ್ರಿಯಾಪದಗಳನ್ನು ೬\par

\begin{kbvderiv}
\vspace{4pt plus 3pt minus 1pt}
\begin{center}\begin{adjustbox}{max width=\textwidth}
{\renewcommand{\arraystretch}{1.22}\setlength{\tabcolsep}{4pt}
\begin{tabular}{l c l}
ಓಡುತ್ತಿದೆ & {\color{kbvgrey}$=$} & ಓಡು + ಉತ್ತ + ಇದೆ \\
ಹಾಡುತ್ತಾಳೆ & {\color{kbvgrey}$=$} & ಹಾಡು + ಉತ್ತ + ಆಳೆ \\
ತಿನ್ನುತ್ತದೆ & {\color{kbvgrey}$=$} & ತಿನ್ನು + ಉತ್ತ + ಅದೆ \\
ಆಡುತ್ತಾನೆ & {\color{kbvgrey}$=$} & ಆಡು + ಉತ್ತ + ಆನೆ \\
\end{tabular}}
\end{adjustbox}\end{center}
\vspace{4pt plus 3pt minus 1pt}
\end{kbvderiv}

\indent ಇಲ್ಲಿ “ಓಡು” ಎಂಬುದು “ಓಡುತ್ತಿದೆ” ಎಂಬ ಕ್ರಿಯಾಪದದ ವ. ರೂಪ; ಇದು ಸ್ವತಂತ್ರವಾದ ಅರ್ಥವನ್ನು ಹೊಂದಿರುವ ಶಬ್ದ. ಕ್ರಿಯೆ ಅರ್ಥವನ್ನು ಕೊಡುವ, ಆದರೆ ಪ್ರತ್ಯಯವನ್ನು ಹೊಂದಿರದ ಈ ಶಃ ಕ್ರಿಯಾಪ್ರಕೃತಿಯಾಗಿರುವುದು. ಕ್ರಿಯಾಪ್ರಕೃತಿಯನ್ನು “ಧಾತು” ಎ ದಾಗಿಯೂ ಕರೆಯುವುದುಂಟು. ಕ್ರಿಯಾಪ್ರಕೃತಿಯ ಮುಂದಿರುವ ಉತ್ತ, ಎಂಬುವು ಪ್ರತ್ಯಯಗಳು. “ಉತ್ತ” ಎಂಬುದನ್ನು “ಕಾಲ ಸೂ ಪ್ರತ್ಯಯ”ವೆಂತಲೂ, “ಇದೆ” ಎಂಬುದನ್ನು ಆಖ್ಯಾತ ಪ್ರತ್ಯಯವೇಂ ಹೇಳುತ್ತೇವೆ.\par

\begin{kbvrule}
\noindent{\kbvsecfont ವ್ಯಾಖ್ಯೆ :}\hspace{0.5em}ಕ್ರಿಯಾಪದದ ಮೂಲ ರೂಪಕ್ಕೆ ಕ್ರಿಯಾಪ್ರಕೃತಿ ೮2 ಧಾತು ಎಂದು ಹೆಸರು. ಸ್ವತಂತ್ರವಾಗಿರುವ ಈ ಧಾತುಗಳೇ ಮ ಧಾತುಗಳು,
\end{kbvrule}
